\section{Conclusion}

This paper presented a comparative study of traditional computer vision and deep learning approaches for herbal plant classification using a dataset constructed from 21 video sequences collected under diverse environmental conditions. A preprocessing pipeline incorporating vegetation filtering, blur detection, image normalization, and data augmentation was developed to improve dataset quality and robustness.

Three classification approaches were evaluated: an ORB-based Bag-of-Visual-Words (BoVW) pipeline with traditional machine learning classifiers, a custom Convolutional Neural Network (CNN), and a MobileNetV2 transfer learning model. Experimental results demonstrated substantial differences in performance among the evaluated methods. The traditional ORB-BoVW-SVM approach achieved moderate classification accuracy, while the custom CNN trained from scratch struggled to generalize effectively due to the limited size of the dataset. In contrast, the MobileNetV2 transfer learning model achieved the highest performance with a test accuracy of 93.23\%, demonstrating the effectiveness of leveraging pre-trained feature representations for plant recognition tasks.

The results indicate that transfer learning is a practical and highly effective solution for image classification problems involving limited training data. By utilizing knowledge learned from large-scale datasets, transfer learning significantly improves classification accuracy while reducing the amount of task-specific data required for training.

Future work could focus on expanding the dataset with additional plant species and environmental conditions to improve model generalization. Further research may also investigate object detection and instance segmentation techniques to localize plant regions automatically within complex scenes using SIFT\cite{sift} or SURF\cite{surf}. Additionally, evaluating more advanced transfer learning architectures such as EfficientNet, ConvNeXt, or Vision Transformers could provide further insights into the effectiveness of modern deep learning approaches for herbal plant recognition.

Overall, the findings of this study demonstrate the feasibility of automated herbal plant classification and provide a reproducible framework for future computer vision applications in agriculture, biodiversity conservation, and traditional medicine.
